
\section{Symbol glossary}\label{app:symbols}
\begin{table}[h]\centering\small
\caption{Notation used throughout.}
\label{tab:symbols}
\begin{tabular}{ll}
\hline
Symbol & Meaning \\
\hline
$s = s_1\cdots s_n$ & peptide sequence, residues from the 20-letter alphabet \\
$L$ & padded sequence length (60 pilot, 150 large, 200 study 14) \\
$K_k(s)$ & set of contiguous $k$-mers of $s$ \\
$J_k$, $\nu_k$ & $k$-mer Jaccard similarity and novelty (Eq.~\eqref{eq:novelty}) \\
$Q(s;p)$ & expected net charge at pH $p$ (Eq.~\eqref{eq:netcharge}) \\
$\mu_H$, $\mu_i(s;w,\delta)$ & Eisenberg hydrophobic moment (Eq.~\eqref{eq:hmoment}) \\
$A$, $A_{\mathrm{norm}}$ & chain-window adjacency and its row normalization \\
$\hat p(s)$ & model posterior $P(\mathrm{AMP}\mid s)$ \\
$t^{*}$ & Bayes decision threshold (Eq.~\eqref{eq:bayesthr}) \\
$A$ (AUROC context) & Mann--Whitney AUROC statistic (Prop.~\ref{prop:mw}) \\
$Q_1, Q_2$ & placement probabilities in the Hanley--McNeil variance \\
$B$ & bootstrap or permutation resample count \\
$\alpha(n)$ & inverse Ackermann factor in union-find complexity \\
$\omega$ & positive-class weight in the weighted BCE \\
$c_{01}, c_{10}$ & false-positive and false-negative costs \\
$\mathrm{BS}$ & Brier score (Eq.~\eqref{eq:brier}) \\
\hline
\end{tabular}
\end{table}
